\clearpage
\thispagestyle{empty}
\pdfbookmark[0]{Notices of the Author’s Own Papers. Edition and sources}{anzeigen-source-notice}
\begin{center}\Large\textbf{Notices of the Author’s Own Papers}\end{center}

AI-translated from German, with AI-produced typesetting. The German source is Carl Friedrich Gauss, \textit{Werke}, volume V (Göttingen, 1867), division title and the notice from the \textit{Göttingische gelehrte Anzeigen} of April 5, 1813, native leaves 283--291. The notice discusses Gauss’s memoir on the attraction of homogeneous ellipsoids.

\noindent\textbf{Translation of Lagrange’s French quotation}\par\smallskip
\noindent Their solutions can only be regarded as masterpieces of analysis, but one may still wish for a more direct and simpler solution, and the continual advances in analysis give reason to hope for one.

\medskip\noindent\textbf{Related exposition}\par\smallskip
\noindent Isaac Todhunter discusses the Latin memoir in \textit{A History of the Mathematical Theories of Attraction and the Figure of the Earth}, vol.~II (1873), §§~1162--1175. He also identifies the historical German notice and its reprint.

\medskip\noindent\href{https://gdz.sub.uni-goettingen.de/id/PPN236006339}{German source (1867)}\quad
\href{https://www.gutenberg.org/ebooks/69832}{Todhunter (1873)}
\clearpage
